Praktické kroky pro zavedení AI do kurzu: jasné výukové cíle, definice automatizovaných částí a kontrolní body pro pilotní fázi \cite{node_2-5}.
Doporučený workflow: \begin{enumerate}
\item Definice cíle/formátu AI (jaké výstupy) \cite{node_2-5}.
\item Archivace promptů/konfigurací a lokální kontrola učitelem \cite{node_2-5}.
\item Pilotní nasazení s měřením (čas, kvalita, spokojenost, výsledky) a iterativní úpravy \cite{node_2-5}.
\end{enumerate}
Volba nástroje a metriky: podporované formáty (LMS, QTI/CSV, multimédia), transparentnost/editovatelnost, bezpečnost/on‑premise; metriky — úspora času, zlepšení výsledků, spokojenost, kvalita materiálů \cite{node_2-5}.
Etika a governance: transparentní komunikace o využití dat, uchovávání verzí promptů pro audit, interní pravidla a profesní rozvoj učitelů; monitorovat rizika (zkreslení, bezpečnost, akademická poctivost) \cite{node_2-5}.